A todo saber corresponde um laço à medida que todo saber aprendido realiza uma expansão do corpo. Todo saber adquirido constitui um novo membro no corpo, que é sensível de algum modo específico a certos elementos até então mudos a tal corpo. Por outro lado, toda relação, todo laço social, se constitui através de um saber que mapeia -- identifica, descreve, relaciona, articula -- o corpo enlaçado no corpo enlaçante de modo a orientar os movimentos, as comunicações entre ambos.